% SPDX-FileCopyrightText: 2025 IObundle
%
% SPDX-License-Identifier: MIT

\documentclass{ug}
\usepackage{xltabular}

% SPDX-FileCopyrightText: 2025 IObundle
%
% SPDX-License-Identifier: MIT

%replace ipcore-name by the name of your ip core (e.g. IOb-Cache) and description by a brief description (e.g. a Configurable Cache)

\title{%
\Huge IOb-System \\
 \vspace*{3cm}
\Large System-on-Chip (SoC) template 
}

\header{IOb-System, System-on-Chip (SoC) template}

\date{\today}
\category{User Guide, \input{\NAME_version.tex}, Build 84f30c0a}

% SPDX-FileCopyrightText: 2025 IObundle
%
% SPDX-License-Identifier: MIT

\definecolor{iob-green}{rgb}{0.0,1.0,0.80}
\definecolor{iob-blue}{rgb}{0.90196,0.94902,1}
\definecolor{dark-blue}{RGB}{155,194,230}
\definecolor{mid-blue}{RGB}{189,215,238}
\definecolor{light-blue}{RGB}{221,235,247}


\begin{document}

\maketitle
\pagenumbering{gobble}

\vspace*{\fill}
User Guide, \input{\NAME_version.tex}, Build 84f30c0a
\hspace*{\fill} \includegraphics[keepaspectratio,scale=.7]{Logo.png}

\cleardoublepage
\pagenumbering{roman}
\setcounter{page}{1}
% SPDX-FileCopyrightText: 2025 IObundle
%
% SPDX-License-Identifier: MIT

\subsection*{Document Version History}

\begin{table}[h]
  \centering
  \begin{tabularx}{\textwidth}{|c|c|c|X|}
    \hline
    \rowcolor{iob-green}
             {\bf Version} & {\bf Date} & {\bf Person} & {\bf Changes from previous version} \\ \hline
             \input ./revhist_desc
  \end{tabularx}
\end{table}

\cleardoublepage
\tableofcontents
\clearpage
\listoftables
\clearpage
\listoffigures
\cleardoublepage
\pagenumbering{arabic}
\setcounter{page}{1}
\section{Introduction}
\label{sec:intro}
% SPDX-FileCopyrightText: 2025 IObundle
%
% SPDX-License-Identifier: MIT

\lipsum[1-1]

% SPDX-FileCopyrightText: 2025 IObundle
%
% SPDX-License-Identifier: MIT

\begin{figure}[!htbp]
  \centerline{\includegraphics[width=14cm]{symb.pdf}}
  \vspace{0cm}
  \caption{IP Core Symbol}
  \label{fig:symbol}
\end{figure}


\subsection{Features}
\label{sec:feat}
% SPDX-FileCopyrightText: 2025 IObundle
%
% SPDX-License-Identifier: MIT

\begin{itemize}
\item 32-bit RISC-V control CPU
\item Support for Integer (I), atomic (A), and multiply/divide extensions (M)
\item Instruction and data caches
\item RS232 interfaces for viewing runtime messages
\item Optional timer peripheral
\item Optional Ethernet peripheral
\item Frequency of operation at 167MHz on FPGA
\item Needs external DDR4 memory controller IP
\end{itemize}


\subsection{Benefits}
\label{sec:benef}
% SPDX-FileCopyrightText: 2025 IObundle
%
% SPDX-License-Identifier: MIT

\begin{itemize}
  \itemsep-0.5em
\item Compact and easy to integrate hardware and software implementation
\item Low power consumption
\end{itemize}


\subsection{Deliverables}
\label{sec:deliv}
% SPDX-FileCopyrightText: 2025 IObundle
%
% SPDX-License-Identifier: MIT

\begin{itemize}
  \itemsep-0.5em
\item FPGA synthesized netlist or Verilog source code, and respective
  synthesis and implementation scripts
\item FPGA verification environment by simulation and emulation
\item Bare-metal software driver and example user software
\item User documentation for easy system integration
\end{itemize}


\ifdefined\SECTIONCLEARPAGE
\clearpage
\fi
\section{Description}

This section gives a detailed description of the IP core. The high-level block
diagram is presented, along with a description of its subblocks. The parameters and
macros that define the core configuration are listed and explained. The
interface signals are enumerated and described; if timing diagrams are needed,
they are shown after the interface signals. Finally the Control and Status
Registers (CSR) are outlined and explained.


\ifdefined\PYPARAMS
\subsection{Python Parameters}
\label{sec:py_params}

The following tables describe the supported \textit{Python Parameters} for the setup of this IP core.
Note that these \textit{Python Parameters} are not used during the build process of this core from this build directory.
They only serve a purpose during the setup process, to configure how the core build directory will be generated.
See the \textit{Python Parameters} section of the \href{https://github.com/IObundle/py2hwsw/blob/main/py2hwsw/py2hwsw_document/document/ug.pdf}{Py2HWSW User Guide} for more details.
\clearpage
\fi

\subsection{Block Diagram}
\label{sec:bdd}
% SPDX-FileCopyrightText: 2025 IObundle
%
% SPDX-License-Identifier: MIT

Figure~\ref{fig:bd} presents a high-level block diagram of the core, followed by
a brief description of each block.

% SPDX-FileCopyrightText: 2025 IObundle
%
% SPDX-License-Identifier: MIT

\begin{figure}[H]
  \centering {\includegraphics[width=\columnwidth]{bd.pdf}}
  \vspace{-0.7cm}
  \caption{High-Level Block Diagram}
  \label{fig:bd}
\end{figure}


\setlist[description]{font=\bfseries\MakeUppercase, labelindent=\parindent, leftmargin=*}

The Verilog modules in the top-level entity of the core are         described in the following tables. The table elements represent         the subblocks in the Block Diagram.

\begin{table}[H]
  \centering
  \begin{tabularx}{\textwidth}{|l|l|X|}

    \hline
    \rowcolor{iob-green}
    {\bf Module} & {\bf Name} & {\bf Description}  \\ \hline \hline

    \input general_operation_subblocks_tab

  \end{tabularx}
  \caption{General operation group}
  \label{general_operation_subblocks_tab:is}
\end{table}
\clearpage


\ifdefined\CONFS
\subsection{Configuration}
\label{sec:ipconfig}

The following tables describe the IP core configuration. The core may be configured using macros or parameters:

\begin{description}
    \item \textbf{'M'} Macro: a Verilog macro or \texttt{define} directive is used to include or exclude code segments, to create core configurations that are valid for all instances of the core.
    \item \textbf{'P'} Parameter: a Verilog parameter is passed to each instance of the core and defines the configuration of that particular instance.
\end{description}

\begin{table}[H]
  \centering
  \begin{tabularx}{\textwidth}{|l|c|c|c|c|X|}

    \hline
    \rowcolor{iob-green}
    {\bf Configuration} & {\bf Type} & {\bf Min} & {\bf Typical} & {\bf Max} & {\bf Description} \\ \hline \hline

    \input general_operation_confs_tab

  \end{tabularx}
  \caption{General operation group}
  \label{general_operation_confs_tab:is}
\end{table}

There are also constants that are used in the core in order to improve the readability
of the code and should not be changed. They are defined as presented in the list below:
\input constants
\clearpage
\fi

\subsection{Interface Signals}
\label{sec:ifsig}
The interface signals of the core are described in the following tables.
Note that the ouput signals are registered in the core, while the input signals are not.
\begin{table}[H]
  \centering
  \begin{tabularx}{\textwidth}{|l|l|r|X|}

    \hline
    \rowcolor{iob-green}
    {\bf Name} & {\bf Direction} & {\bf Width} & {\bf Description}  \\ \hline \hline

    \input clk_en_rst_s_if_tab

  \end{tabularx}
  \caption{Clock, clock enable and reset}
  \label{clk_en_rst_s_if_tab:is}
\end{table}

\begin{table}[H]
  \centering
  \begin{tabularx}{\textwidth}{|l|l|r|X|}

    \hline
    \rowcolor{iob-green}
    {\bf Name} & {\bf Direction} & {\bf Width} & {\bf Description}  \\ \hline \hline

    \input iob_csrs_cbus_s_if_tab

  \end{tabularx}
  \caption{Control and Status Registers interface (auto-generated)}
  \label{iob_csrs_cbus_s_if_tab:is}
\end{table}
\clearpage

%timing diagrams
\ifdefined\TD
\subsection{Timing Diagrams}
\label{sec:td}
\input{td}
\fi

%software components
\ifdefined\CSRS
\subsection{Control and Status Registers}
\label{sec:csrs}

The software accessible registers of the core are described in the following
tables. The tables give information on the name, read/write capability, address, width in bits, and a textual description.
\subsubsection{ Configuration}
\begin{table}[H]
  \centering
  \begin{tabularx}{\textwidth}{|l|c|c|c|c|X|}

    \hline
    \rowcolor{iob-green}
    {\bf Name} & {\bf R/W} & {\bf Addr} & {\bf Width} & {\bf Default} & {\bf Description} \\ \hline

    \input _uart_csrs_tab

  \end{tabularx}
  \caption{UART software accessible registers}
  \label{_uart_csrs_tab:is}
\end{table}

\begin{table}[H]
  \centering
  \begin{tabularx}{\textwidth}{|l|c|c|c|c|X|}

    \hline
    \rowcolor{iob-green}
    {\bf Name} & {\bf R/W} & {\bf Addr} & {\bf Width} & {\bf Default} & {\bf Description} \\ \hline

    \input _general_operation_csrs_tab

  \end{tabularx}
  \caption{General Registers.}
  \label{_general_operation_csrs_tab:is}
\end{table}
\clearpage
\fi

%memories
\ifdefined\MEMS
\subsection{Memories}
\label{sec:mems}

    The memories of the core are described in the following table.
    The tables give information on the name, type, address and data width in bits, and initialization file name.

    \begin{table}[H]
      \centering
      \begin{tabularx}{\textwidth}{|l|c|c|c|X|}
        \hline
        \rowcolor{iob-green}
        {\bf Name} & {\bf Type} & {\bf (Word-)Addr Width} & {\bf Data Width} & {\bf Init file} \\ \hline
        \input mems_tab
      \end{tabularx}
      \caption{Table of memories of the core}
      \label{mems_tab:is}
    \end{table}
\clearpage
\fi

\ifdefined\SECTIONCLEARPAGE
\clearpage
\fi
\section{Usage}

\subsection{Instantiation}
\label{sec:inst}
% SPDX-FileCopyrightText: 2025 IObundle
%
% SPDX-License-Identifier: MIT

Figure~\ref{fig:inst} illustrates how to instantiate the IP core and, if applicable, the
required external subblocks.

\begin{figure}[!htbp]
    \centerline{\includegraphics[width=14cm]{inst.pdf}}
    \vspace{0cm}\caption{Core Instance and Required Surrounding Blocks}
    \label{fig:inst}
\end{figure}

% SPDX-FileCopyrightText: 2025 IObundle
%
% SPDX-License-Identifier: MIT

IOb-System by default contains a RS232 interface that should be connected to an external UART (e.g. a USB-to-serial converter).

The AXI bus for the external memory is optional according to the Python Parameter configuration of the system. If the system has this port, then it must be connected.

The clock, clock enable, and reset ports can be connected to the desired clock and reset generator.

The IOb-System may also contain memory ports for the boot ROM and RAM, but these memories are instantiated automatically in the IOb-System Memory Wrapper.




\subsection{Simulation}
\label{sec:tbbd}
% SPDX-FileCopyrightText: 2025 IObundle
%
% SPDX-License-Identifier: MIT

The provided testbench uses the core instance described in
Section~\ref{sec:inst}. A high-level block diagram of the testbench is shown in
Figure~\ref{fig:tbbd}. The testbench is organized in a modular fashion, with
each test described in a separate file. The test suite consists of all the test
case files to make adding, modifying, or removing tests easy.

\begin{figure}[!htbp]
    \centerline{\includegraphics[width=14cm]{tbbd.pdf}}
    \vspace{0cm}\caption{Testbench Block Diagram}
    \label{fig:tbbd}
\end{figure}

% SPDX-FileCopyrightText: 2025 IObundle
%
% SPDX-License-Identifier: MIT

The IOb-System simulation wrapper contains its own UART core that connects to the system's RS232 port.

Depending on the Python Parameter configuration of the system, the simulation wrapper may contain an external memory.

The clock, clock enable, and reset ports are connected to the testbench clock and reset generator.


\subsubsection*{System-level Simulation}

Upon request, simulation files to run the core embedded in a RISC-V system can
be provided. The core is exercised in various modes by the RISC-V processor,
using a bare-metal software program written in the C programming language.


\ifdefined\ASICSYNTH
\subsection{ASIC Synthesis}
\label{sec:synth}
% SPDX-FileCopyrightText: 2025 IObundle
%
% SPDX-License-Identifier: MIT

A {\tt .tcl} script is provided for the Cadence Genus synthesis tool. The script
reads the technology files, compiles and elaborates the design, and proceeds to
synthesize it. Timing constraints file(s) are provided.

After synthesis, reports on silicon area usage, power consumption, and timing
closure are generated. A post-synthesis Verilog file is created, which can be
used in post-synthesis simulation. The IP core has been synthesized without the
FIFO memories, which are treated as black boxes. RAM subblocks are not included as
they are technology dependent. The user is free to generate and add the required
memories, or request this service from the IP core provider.

\fi

\ifdefined\FPGACOMP
\subsection{FPGA Compilation}
\label{sec:fpga}
% SPDX-FileCopyrightText: 2025 IObundle
%
% SPDX-License-Identifier: MIT

The Quartus and Vivado FPGA compilation toolchains are supported via {\tt .tcl}
scripts invoked by a Makefile. The script compiles and elaborates the design for
a given set of target FPGA boards. Additional boards can be supported by
following the flow provided.

The core is instantiated within a Verilog wrapper, which depends on the
specifics of the FPGA device and the FPGA board. Timing constraints file(s) are
provided.

After compilation, reports on FPGA resource usage, power consumption, and timing
closure are generated. A post-synthesis Verilog file is created, which can be
used in post-synthesis simulation.

% SPDX-FileCopyrightText: 2025 IObundle
%
% SPDX-License-Identifier: MIT

Add here any specific ip core description.


\subsubsection*{System-level FPGA Run}

Upon request, files to run the core on the supported FPGA boards can be
provided. The core is embedded in a RISC-V system and exercised in various
modes, using a bare-metal software program written in the C programming
language.

\fi

\ifdefined\RESULTS
\ifdefined\SECTIONCLEARPAGE
\clearpage
\fi
\section{Implementation Results}
\label{sec:results}
% SPDX-FileCopyrightText: 2025 IObundle
%
% SPDX-License-Identifier: MIT

The results obtained use all the default user configurable parameters, which affect the RAM usage results.

\fi

\ifdefined\CUSTOM
\input{custom}
\fi

\end{document}
